\chapter{多项式}
\orig{1--5}
这一章，我们主要是讨论数域 $P$ 上的一元多项式，并举出有关的反例。关于整除、最大公因式、互素、不可约、$k$ 重因式以及本原多项式这一系列的概念，都为大家所熟悉，就不一一叙述了。

下面就举一些相关的反例。

\begin{enumerate}
\item 如果 $f(x)\mid g_i(x)$，$i=1,2,\ldots,n$，那么 $f(x)$ 就能整除 $g_1(x),g_2(x),\ldots,g_n(x)$ 的组合，即
\[
 f(x)\mid\bigl(u_1(x)g_1(x)+u_2(x)g_2(x)+\cdots+u_n(x)g_n(x)\bigr).
\]
\textbf{反之不真。} 即 $f(x)$ 能整除 $g_1(x),g_2(x),\ldots,g_n(x)$ 的组合，未必有 $f(x)\mid g_i(x)$ 对每个 $i$ 都成立。

\textbf{例}
\[
 f(x)=3x-2,
 \qquad g_1(x)=x^2+1,
 \qquad g_2(x)=2x+3,
\]
\[
 u_1(x)=-2,
 \qquad u_2(x)=x.
\]
显然
\[
 u_1(x)g_1(x)+u_2(x)g_2(x)
 =-2(x^2+1)+x(2x+3)=3x-2=f(x),
\]
故
\[
 f(x)\mid\bigl(u_1(x)g_1(x)+u_2(x)g_2(x)\bigr),
\]
但 $f(x)\nmid g_1(x)$。

\item $P[x]$ 中的多项式 $f(x)$ 与 $g(x)$ 的最大公因式是 $d(x)$，那么有 $P[x]$ 中的多项式 $u(x)$ 与 $v(x)$，使
\[
 d(x)=u(x)f(x)+v(x)g(x).
\]
\jiaokan{此处排印为形似 $d(x)=u(x)(f(x)+v(x)g(x)$ 的式子，多出一个左括号且破坏了 Bézout 恒等式。依上下文及随后各例，应为 $d(x)=u(x)f(x)+v(x)g(x)$。}
但：
\begin{enumerate}[label=\roman*),leftmargin=2.4em,itemsep=0.35em]
\item 反之不真。即上式成立，$d(x)$ 未必是 $f(x)$ 与 $g(x)$ 的最大公因式；
\item 满足上式的 $u(x)$ 与 $v(x)$ 不是唯一的。
\end{enumerate}

\textbf{例 i)} 令
\[
 f(x)=x,\qquad g(x)=x+1.
\]
有
\[
 x(x+2)+(x+1)(x-1)=2x^2+2x-1,
\]
其中
\[
 u(x)=x+2,\qquad v(x)=x-1.
\]
而 $2x^2+2x-1$ 显然不是 $f(x)$ 与 $g(x)$ 的最大公因式。

我们说，当
\[
 d(x)=u(x)f(x)+v(x)g(x)
\]
而 $d(x)$ 是 $f(x)$ 与 $g(x)$ 的一个公因式时，$d(x)$ 一定是 $f(x)$ 与 $g(x)$ 的一个最大公因式。

\textbf{例 ii)} 令
\[
 f(x)=x^2-1,\qquad g(x)=1,
\]
则 $d(x)=1$。例如下列三组 $u,v$ 都满足 Bézout 恒等式：
\[
\begin{cases}
 u(x)=-1,\\
 v(x)=x^2,
\end{cases}
\qquad
\begin{cases}
 u(x)=0,\\
 v(x)=1,
\end{cases}
\qquad
\begin{cases}
 u(x)=-2,\\
 v(x)=2x^2-1.
\end{cases}
\]

\item 多项式 $f_1(x),f_2(x),\ldots,f_n(x)$（$n>2$）互素时，并不一定两两互素。

\textbf{例 i)}
\[
\begin{aligned}
 f_1(x)&=x^2-3x+2,\\
 f_2(x)&=x^2-5x+6,\\
 f_3(x)&=x^2-4x+3.
\end{aligned}
\]
它们整体互素，但
\[
 (f_1(x),f_2(x))=x-2.
\]

\textbf{例 ii)}
\[
\begin{aligned}
 f_1(x)&=x^3-7x^2+7x+15,\\
 f_2(x)&=x^2-x-20,\\
 f_3(x)&=x^3+x^2-12x.
\end{aligned}
\]
则
\[
 (f_1,f_2,f_3)=1,
\]
而
\[
 (f_1,f_2)=x-5,
 \qquad (f_1,f_3)=x-3,
 \qquad (f_2,f_3)=x+4.
\]

\item 不可约多项式 $p(x)\mid f(x)g(x)$，则有 $p(x)\mid f(x)$ 或 $p(x)\mid g(x)$。

我们说，$p(x)$ 是不可约多项式的限制是必要的，否则即可举出以下反例：
\[
 p(x)=(x-1)^2,\qquad f(x)=x-1,\qquad g(x)=x-1.
\]
显然 $p(x)\mid f(x)g(x)$，但 $p(x)\nmid f(x)$ 且 $p(x)\nmid g(x)$。

\item 若不可约多项式 $p(x)$ 是 $f(x)$ 的 $k$ 重因式（$k\ge 1$），则 $p(x)$ 是 $f'(x)$ 的 $k-1$ 重因式。

\textbf{反之不真。} 举例如下：
\[
 f(x)=x^3-3x^2+3x-3,
\]
则
\[
 f'(x)=3x^2-6x+3=3(x-1)^2.
\]
$x-1$ 是 $f'(x)$ 的 2 重因式，但 $x-1$ 不是 $f(x)$ 的 3 重因式；其实，它根本不是 $f(x)$ 的因式。

\item 本原多项式不一定是不可约的。

\textbf{例}
\[
 x^2+3x+2
\]
是本原的，可是
\[
 x^2+3x+2=(x+1)(x+2).
\]

\item 设 $f(x),g(x)$ 是整系数多项式，且 $g(x)$ 是本原的；若
\[
 f(x)=g(x)h(x),
\]
其中 $h(x)$ 是有理系数多项式，则 $h(x)$ 一定是整系数的。

我们说，$g(x)$ 限制为本原的条件不可缺少，否则就有
\[
 f(x)=x^2+x,\qquad g(x)=2x+2,
\]
而
\[
 f(x)=g(x)h(x),\qquad h(x)=\frac12x.
\]

\item Eisenstein 判别法告诉我们：当
\[
 f(x)=a_nx^n+a_{n-1}x^{n-1}+\cdots+a_0
\]
是一个整系数多项式，存在一个素数 $p$ 使得
\[
 p\nmid a_n,
\]
\[
 p\mid a_i,\qquad i=0,1,\ldots,n-1,
\]
且
\[
 p^2\nmid a_0,
\]
那么 $f(x)$ 在有理数域上是不可约的。

可是当找不到这样的素数 $p$ 时，我们不能判定其是否可约。

\textbf{例}
\[
 f_1(x)=x^2+3x+2,
 \qquad
 f_2(x)=x^2+1.
\]
对 $f_1(x)$ 及 $f_2(x)$ 来说，都找不到满足判别法条件的素数 $p$；但 $f_1(x)$ 可约，而 $f_2(x)$ 不可约。

\item $f(x)$ 在有理数域上可约，它不一定有有理根。

\textbf{例}
\[
 f(x)=(x^2+1)^2.
\]
\end{enumerate}

这一章的最后，给出关于多项式函数的反例。

\textbf{定义} 如果在多项式 $f(x)$ 与 $g(x)$ 中，同次项的系数全相等，那么 $f(x)$ 与 $g(x)$ 就叫做相等，记为
\[
 f(x)=g(x).
\]

\textbf{定义} 由一个多项式定义的函数称为域 $P$ 上的\textbf{多项式函数}。即设
\[
 f(x)=a_0x^n+a_1x^{n-1}+\cdots+a_n
\]
是 $P[x]$ 中的多项式，$\alpha$ 是 $P$ 中的数，在 $f(x)$ 的表达式中用 $\alpha$ 代 $x$ 所得的数
\[
 a_0\alpha^n+a_1\alpha^{n-1}+\cdots+a_n
\]
称为 $f(x)$ 当 $x=\alpha$ 时的值，记为 $f(\alpha)$。这样一来，多项式 $f(x)$ 就定义了一个域 $P$ 上的函数。

我们知道，两个多项式相等必然导出两个多项式函数相等；反之，两个多项式函数相等是否必然导出两个多项式相等呢？

当 $P$ 是无限域时，回答是肯定的；当 $P$ 是有限域时，回答则未必。\jiaokan{此处写作“当 $P$ 是数域时……当 $P$ 是有穷多个元素时……”。按旧式术语，“数域”通常指复数域的子域，因而必为无限域；随后又使用有限域 $\mathbb Z_2$。为使逻辑条件直接对应结论，此处改为“无限域/有限域”。}

\textbf{例} $\mathbb Z_2$ 是以 2 为模的剩余类域。考察 $\mathbb Z_2$ 上多项式
\[
 f(x)=x+1,
 \qquad
 g(x)=x^2+1.
\]
有
\[
 f(0)=g(0)=1,
 \qquad
 f(1)=g(1)=0.
\]
所以，这里的两个多项式函数 $f(x)$ 与 $g(x)$ 相等，但是这里的两个多项式 $f(x)$ 与 $g(x)$ 不等。
